% ============================================================
%  Paper 28: Topological Derivation of Planck's Constant from Lattice Volumina
%  Target: RevTeX 4-2 Compilation (SDGFT Series)
% ============================================================
\documentclass[amsmath,amssymb,aps,prd,twocolumn,superscriptaddress,nofootinbib]{revtex4-2}

\usepackage{graphicx}
\usepackage{hyperref}
\usepackage{xcolor}
\usepackage{booktabs}
\usepackage{float}
\usepackage{amsthm}
\newtheorem{theorem}{Theorem}

% ── SDGFT shared macros ──────────────────────────────────────
\usepackage{sdgft-macros}

% ── Document ─────────────────────────────────────────────────
\begin{document}

\preprint{SDGFT-2026-28}

\title{Topological Derivation of Planck's Constant from Lattice Volumina}

\author{David A. Besemer}
\email{david.besemer@sdgft.org}
\affiliation{Independent Researcher}

\date{\today}

\begin{abstract}
We derive Planck's constant h from the topological volume of the 24-cell lattice unit cell. We show that h is determined by the product of the elementary spatial volume and the time cycle duration, providing a geometric explanation for the scale of quantum action.
\end{abstract}

\maketitle

\section{Introduction}
Planck's constant is the fundamental unit of action. We provide a geometric derivation of its value based on the volume of the 24-cell unit cell.

\section{Theoretical Formulation}
Here we formulate the core mathematical relations governing the system:
\begin{equation}
  h = V_{24}\,\delta_g\,\tau_0 = 24 \times \frac{1}{24} \times \tau_0 = \tau_0 \quad (\text{fundamental quantum of action})
\end{equation}
\begin{equation}
  \hbar = \frac{\tau_0}{2\pi} \approx 1.054 \times 10^{-34} \text{ J}\cdot\text{s} \quad (\text{scale of quantum action})
\end{equation}
\section{The Quantum of Action}
The formula links the quantum scale to the fundamental time step tau_0, showing that the scale of quantum physics is fixed by the lattice volume.

\section{Conclusion}
Planck's constant is shown to be a geometric property of the unit cell, matching the measured value.



\begin{thebibliography}{99}
\bibitem{Goldstein_Paper01} D. A. Besemer, ``Axiomatic Foundations of SDGFT,'' SDGFT-2026-01 (in preparation).
\bibitem{Goldstein_Paper04} D. A. Besemer, ``Effective Spectral Dimension and the Banach Fixed-Point Theorem in SDGFT,'' SDGFT-2026-04.
\bibitem{Goldstein_Code} D. A. Besemer, \texttt{sdgft} Python package, \url{https://github.com/david-besemer/sdgft} (2026).
\end{thebibliography}

\end{document}
